\documentclass{article}

\usepackage[utf8]{inputenc}
\usepackage[T1]{fontenc}
\usepackage{hyperref}
\usepackage{url}
\usepackage{booktabs}
\usepackage{amsfonts}
\usepackage{amsmath}
\usepackage{microtype}
\usepackage{graphicx}

% Results are generated from a reproducible experiment, not typed in by hand:
%   make -C papers results        (writes papers/generated/*.tex)
% Every \mc... macro below expands to a number produced by that run.
% AUTO-GENERATED by papers/generate_results.py -- do not edit by hand.
% Regenerate with: make -C papers results
\newcommand{\mcCalibratedAccuracy}{0.6987}
\newcommand{\mcCalibratedEce}{0.016}
\newcommand{\mcCalibratedMce}{0.0336}
\newcommand{\mcCalibratedBrier}{0.1909}
\newcommand{\mcCalibratedAuroc}{0.692}
\newcommand{\mcOverconfidentAccuracy}{0.6987}
\newcommand{\mcOverconfidentEce}{0.1231}
\newcommand{\mcOverconfidentMce}{0.1781}
\newcommand{\mcOverconfidentBrier}{0.2088}
\newcommand{\mcOverconfidentAuroc}{0.6929}
\newcommand{\mcUnderconfidentAccuracy}{0.6987}
\newcommand{\mcUnderconfidentEce}{0.2769}
\newcommand{\mcUnderconfidentMce}{0.3539}
\newcommand{\mcUnderconfidentBrier}{0.2704}
\newcommand{\mcUnderconfidentAuroc}{0.6929}
\newcommand{\mcRandomAccuracy}{0.6987}
\newcommand{\mcRandomEce}{0.2879}
\newcommand{\mcRandomMce}{0.6911}
\newcommand{\mcRandomBrier}{0.3368}
\newcommand{\mcRandomAuroc}{0.506}


\title{Knowing What You Know: A Reproducible Harness for\\
Metacognitive Calibration and Introspective Self-Prediction}

\author{
  Hidden Layer Lab \\
  \texttt{contact@hiddenlayer.ai}
}

\date{}

\begin{document}

\maketitle

\begin{abstract}
A trustworthy model should know what it knows: its stated confidence should track
whether it is actually correct, and---more sharply---it should be able to predict its
own outputs. We present a small, provider-agnostic harness for measuring two facets of
this self-knowledge: \emph{metacognitive calibration} (do confidences match
correctness?) and \emph{introspective self-prediction} (can a model forecast its own
answer?). The harness computes Brier score, Expected and Maximum Calibration Error
(ECE/MCE), and a confidence-vs-correctness AUROC, and runs fully offline through a
deterministic simulated subject so that results are reproducible on any machine without
API access. As a validation, we sweep four calibration profiles and recover the expected
ordering: a calibrated subject attains ECE $=\mcCalibratedEce$ while an overconfident
variant inflates to ECE $=\mcOverconfidentEce$ at identical task accuracy, and a subject
whose confidences carry no signal collapses to chance discrimination
(AUROC $=\mcRandomAuroc$). Every number in this paper is generated directly from the
experiment, illustrating a reproducible path from hypothesis to publication.
\end{abstract}

\section{Introduction}
Calibration is a classical desideratum for probabilistic forecasts \cite{brier1950}, and
has re-emerged as a central concern for modern neural networks, which are often
miscalibrated despite high accuracy \cite{guo2017}. For language models the question is
sharpened by \emph{introspection}: beyond reporting a well-calibrated confidence, can a
model anticipate its own behavior \cite{anthropic2025}? These are distinct,
independently measurable properties, and both are candidate alignment signals---a system
that reliably flags ``I don't know'' is safer to defer to a human.

We contribute a minimal harness that (i) elicits an answer and a verbalized confidence,
(ii) scores correctness deterministically, and (iii) reports a standard battery of
calibration and discrimination metrics. The harness is offline-first: a deterministic
simulated subject with tunable accuracy and calibration lets us validate the metrics and
run continuous-integration checks with no API keys.

\section{Method}
\paragraph{Metrics.} For predictions with stated confidence $c_i \in [0,1]$ and outcomes
$o_i \in \{0,1\}$, we report the Brier score $\frac{1}{N}\sum_i (c_i - o_i)^2$; Expected
Calibration Error, the count-weighted mean gap $\sum_b \frac{n_b}{N}\,|\,\text{conf}_b -
\text{acc}_b|$ over confidence bins; Maximum Calibration Error, the worst such gap; and
the area under the ROC curve using confidence to discriminate correct from incorrect
answers. Calibration (ECE/MCE/Brier) and discrimination (AUROC) are orthogonal: a model
can rank its right answers above its wrong ones yet still be systematically over- or
under-confident.

\paragraph{Simulated subjects.} To validate the metrics we use a deterministic subject
that, per item, draws a latent correctness probability $p$ near a target accuracy, samples
correctness $\sim\text{Bernoulli}(p)$, and emits a stated confidence as a function of $p$:
\emph{calibrated} ($\text{conf}=p$), \emph{overconfident} (inflated toward $1$),
\emph{underconfident} (deflated), or \emph{random} (independent of $p$). All randomness is
seeded per item, so the experiment is bit-for-bit reproducible.

\section{Results}
We evaluate all four profiles at target accuracy $0.7$ on $1{,}500$ deterministic
arithmetic items. Table~\ref{tab:mc-calibration} reports the full battery; all values are
produced by the experiment and injected via generated macros.

% AUTO-GENERATED by papers/generate_results.py -- do not edit by hand.
% Regenerate with: make -C papers results
\begin{table}[t]
\centering
\begin{tabular}{lrrrrr}
\toprule
arm & accuracy & ece & mce & brier & auroc \\
\midrule
calibrated & 0.6987 & 0.016 & 0.0336 & 0.1909 & 0.692 \\
overconfident & 0.6987 & 0.1231 & 0.1781 & 0.2088 & 0.6929 \\
underconfident & 0.6987 & 0.2769 & 0.3539 & 0.2704 & 0.6929 \\
random & 0.6987 & 0.2879 & 0.6911 & 0.3368 & 0.506 \\
\bottomrule
\end{tabular}
\caption{Calibration of four simulated subjects on 1{,}500 deterministic factual items (offline \texttt{sim} provider, seed 7). ECE/MCE/Brier lower is better; AUROC measures whether confidence ranks correct above incorrect answers.}
\label{tab:mc-calibration}
\end{table}


The recovered ordering matches theory. At identical accuracy ($\mcCalibratedAccuracy$),
the calibrated subject is well-calibrated (ECE $=\mcCalibratedEce$), the overconfident
subject inflates confidence (ECE $=\mcOverconfidentEce$, MCE $=\mcOverconfidentMce$), and
the underconfident subject deflates it (ECE $=\mcUnderconfidentEce$). Crucially, the
\emph{random} profile---whose confidences carry no signal---has chance discrimination
(AUROC $=\mcRandomAuroc$) while the other three share near-identical AUROC
($\approx \mcCalibratedAuroc$), empirically separating calibration from discrimination.

\IfFileExists{figures/metacognition_reliability.png}{%
\begin{figure}[t]
\centering
\includegraphics[width=0.55\textwidth]{figures/metacognition_reliability.png}
\caption{Reliability diagram for the overconfident subject: stated confidence sits
systematically above empirical accuracy.}
\label{fig:reliability}
\end{figure}}{}

\section{Discussion}
The harness turns ``does the model know what it knows?'' into a small set of numbers that
are cheap to compute and hard to game (Brier and the log-loss family are proper scoring
rules). Because the same code path runs against any provider, a researcher can drop in a
real model and obtain comparable metrics; the simulated subjects serve as analysis-time
ground truth and as regression fixtures. The companion self-prediction task extends this
to introspection proper: a subject first \emph{predicts} its answer and then produces it,
and we report self-prediction accuracy separately from task accuracy, since genuine
self-knowledge means forecasting one's own answer whether or not it is correct.

\paragraph{Reproducibility.} This paper is itself an artifact of the pipeline it
describes. The numbers above are not transcribed; running \texttt{make -C papers results}
re-executes the experiment and rewrites the macros and table that this document
\texttt{\textbackslash input}s.

\section{Limitations}
The validation uses simulated subjects, so it tests the \emph{measurement} apparatus, not
any real model's metacognition. Verbalized confidence is only one elicitation method;
token-probability and consistency-based confidences may behave differently. Calibration
estimates are bin- and sample-size-sensitive, which is why we validate in the large-$N$
regime.

\section{Conclusion}
We presented a compact, reproducible harness for metacognitive calibration and
introspective self-prediction, validated it by recovering the expected behavior of four
calibration profiles, and demonstrated a numbers-from-runs path to publication. The same
infrastructure applies directly to frontier models as an alignment-relevant probe of
self-knowledge.

\bibliographystyle{plain}
\bibliography{references}

\end{document}
